
\subsection*{1.2.}
Soit maintenant $X$ un schéma quelconque et considérons le groupe $K^{\bullet}(X)$ des classes de $\underline{O}_X$-modules localement libres (par la suite sous-entendu : de type fini). C'est une pré-$\lambda$-algèbre augmentée sur l'anneau $H^0(X,\mathbb Z)$ (V 3.9.1). L'application $f\rightsquigarrow f\cdot1\in K^{\bullet}(X)$ permet d'identifier $H^0(X,\mathbb Z)$ par la suite avec un sous-anneau de $K^{\bullet}(X)$. Si $\underline{E}$ est un $\underline{O}_X$-module localement libre, on désignera par $\cl^{\bullet}(\underline{E})$ sa classe dans $K^{\bullet}(X)$. On définit les opérateurs $\lambda^k$ ($k\in\mathbb N$) par
\begin{equation}\tag{1.2.1}
\lambda^k(\cl^{\bullet}(\underline{E}))=\cl^{\bullet}(\Lambda_{\underline{O}}^k\underline{E})
\end{equation}
et par
\begin{equation}\tag{1.2.2}
\lambda^k(x+y)=\sum_{\substack{i+j=k\\ i,j\in\mathbb N}}\lambda^i(x)\lambda^j(y).
\end{equation}
(V 2.2). Rappelons que dans la $\lambda$-filtration $(\Filt^i(X))_{i\in\mathbb Z}$ de $K^{\bullet}(X)$ on a $\Filt^i(X)=K^{\bullet}(X)$ pour $i\leq0$, et que pour $i\geq1$, $\Filt^i(X)$ est l'idéal engendré par les produits
\begin{equation}\tag{1.2.3}
P=\gamma^{i_1}(x_1)\cdots\gamma^{i_k}(x_k)
\end{equation}
avec $1\leq i_1,\ldots,i_k$ ; $i_1+\cdots+i_k\geq i$ et $x_1,\ldots,x_k\in I=\Ker(\varepsilon)$ ($\varepsilon$ étant l'homomorphisme d'augmentation, cf. Exposé V 3.10). On a démontré (V 3.10) que $\Filt^i(X)$ est même le sous-$H^0(X,\mathbb Z)$-module (donc le sous-groupe abélien, si $X$ est connexe) engendré par les produits (1.2.3), et où on peut de plus supposer que les éléments $x_p$ ($p=1,\ldots,k$) sont de la forme
\begin{equation}\tag{1.2.4}
x_p=\cl^{\bullet}(\underline{E}_p)-\underline{n}_p,
\end{equation}
$\underline{E}_p$ étant un $\underline{O}_X$-module localement libre de rang $\underline{n}_p\in H^0(X,\mathbb Z)$.

